\cotefile{156-8}
\scanpage{10}
\mspage{\unsure{1bis|bis}}
% archive sheet number: 10
\margin{27 juin}
% slanted

Avant de passer aux figures, il faudrait donner quelques compléments aux \unsure{$M$} liminaires

\deleted{\uline{Prop} Si $Y\overset{\circ}{\ll}X$ et $Y\overset{\unsure{\circ}}{\ll}X'$ alors $X\unsure{\gtrless|\bowtie}$}

\deleted{\uline{Prop} Supposons $X\gtrless Y$. Si $Z\overset{\circ}{\ll}X$, alors \unsure{$\ll$}}
\[
\deleted{Z\ll\unsure{Y}\ \Longleftrightarrow\ \unsure{X}\ll\unsure{Y}}
\]
\deleted{\uline{Dém} $\Longleftarrow$ \unsure{trivial}, prouvons $\Longrightarrow$. Par $\unsure{\mathrm{Mog}|\mathrm{Mag}}\,2$, $\exists\,Y'$ \unsure{avec}}
\[
\deleted{Z\overset{\circ}{\ll}Y'\leq Y}
\]
\noindent\fbox{\parbox{.94\linewidth}{\uline{Prop 1} \unsure{Si} $X'\leq X$, $Y'\leq Y$, alors $X\gtrless Y\Longrightarrow X'\gtrless Y'$}}

\uline{Prop 1.} a) $X\gtrless Y$ est symétrique et réflexive

b) Si $X'\leq X$, $Y'\leq Y$, alors $X\gtrless Y\Longrightarrow X'\gtrless Y'$

\uuline{Dém.} a) \deleted{\unsure{symétrie}} trivial par \uuline{Mog 3 a)}.

b) Soit $X''\in\widetilde X'\setminus\widetilde X'\cap\widetilde Y'$, $Y''\in\widetilde Y'\setminus\widetilde X'\cap\widetilde Y'$. \inserted{A prouver $X''\mathrel{|\unsure{\circ|0}|}Y''$.} Si $X'',Y''\notin\widetilde X\cap\widetilde Y$, alors $X''\mathrel{|\circ|}Y''$ par déf de $X\gtrless Y$. Si $X''\in\widetilde X\cap\widetilde Y$, on a $X''\in\widetilde Y$, $Y''\in\widetilde Y'\subset\widetilde Y$, d'autre part évidemment $X''\neq Y''$, donc $X''\mathrel{|\circ|}Y''$ par \uuline{Mog 4}.

\margin{\uline{Prop. 2.} Si $X\gtrless Y$, alors $X\mathrel{|\circ|}Y\Longleftrightarrow X\neq Y$. \uuline{dém} $\Longrightarrow$ par Mog 3 a). Inv. supp. $X\neq Y$, \deleted{\unsure{et}} si $X\nleq Y$ et $Y\nleq X$, i.e. $X\notin\widetilde X\cap\widetilde Y$, $Y\notin\widetilde X\cap\widetilde Y$, on a $X\mathrel{|\circ|}Y$ par déf. de $X\gtrless Y$. Si $X\leq Y$, ou $Y\leq X$, on applique \uline{Mog 4}}
% margin note written sideways

\deleted{\uline{Prop \unsure{2}}}\inserted{\uline{Cor 1}} Soient $X,Y$ avec $X\gtrless Y$. Soit $Z$ \unsure{avec} $Z\overset{\circ}{\ll}X$.

Alors
\[
Z\ll Y\ \Longleftrightarrow\ \deleted{Z}X\leq Y
\]

Énonçons tout de suite corollaire

\uline{Corollaire 2} Si $Z\overset{\circ}{\ll}X$, $Z\overset{\circ}{\ll}Y$, et si $X\gtrless Y$, alors $X=Y$.

En effet, on a $X\ll Y$ et $Y\ll X$, donc $X=Y$ par \uline{Mog 1 a}.

Prouvons \deleted{\uline{Prop 2}}\inserted{cor 1}. $\Longleftarrow$ évident, prouvons
\[
Z\overset{\circ}{\ll}X,\ Z\ll Y\ \Longrightarrow\ X\ll Y\qquad(\text{si }X\gtrless Y).
\]
\deleted{\unsure{Soit}} Par \uuline{Mog 2}, $\exists\,Y'$ \unsure{avec} $Z\overset{\circ}{\ll}Y'\leq Y$, par prop $1$ on a $X\gtrless Y'$, il suffit de prouver $X=Y'$ (d'où $\underset{\inserted{\substack{\Vert\\Y'}}}{X}\leq Y$), i.e. on est ramené au cas du cor. Mais si on avait $X\neq Y$

\[
\boxed{Z\overset{\circ}{\ll}X,\ Z\overset{\circ}{\ll}Y,\ X\gtrless Y\ \Longrightarrow\ X=Y.}\quad\deleted{\unsure{Cor}}
\]

\deleted{\unsure{Prouvons} $Z\unsure{\notin|\in}\widetilde X\cap\widetilde Y$, \unsure{on a}}
\margin{\deleted{$Z\in\widetilde X\setminus\widetilde X\cap\widetilde Y$ \unsure{?}}}

\deleted{\deleted{$Z\in$} Il suffit de prouver $X\neq Y$ (car on a de même $Y\leq X$). on aurait $\widetilde X\subset\widetilde Y$, \deleted{sinon} \unsure{on|ou} aurait $\widetilde X\in\widetilde X\cap\widetilde Y$. Sinon, on aurait $X\mathrel{|\circ|}Y$ (cf lemme), \deleted{d'où} et comme $Z\overset{\circ}{\ll}X$, $Z\overset{\circ}{\ll}Y$ on conclurait $Z\mathrel{|\circ|}Z$ par Mog 3 b), absurde par Mog 3 a). \unsure{Ainsi}}

on aurait par prop 2 $X\mathrel{|\circ|}Y$, donc $Z\mathrel{|\circ|}Z$ par Mog 3 b), absurde par Mog 3 a).
